Use the Tate convention in which the cyclotomic character has
Hodge--Tate weight $-1$. Apply \cite[Theorem 1.0.4]{pan2022} at
$\ell=5$; there is no residual-image proviso at this prime.
It gives a cuspidal eigenform up to character twist. The twisting
character is a de Rham subrepresentation of the Hom representation
and is unramified outside finitely many primes. Local algebraicity
\cite[Appendix B, pp.~32--33]{conrad-local-characters} and class field
theory over $\mathbb Q$ make it a finite character times an integral
power of the cyclotomic character: removing its integral local
Hodge--Tate weight leaves finite inertia at five and at every other
prime, hence a finite ray-class character. Comparing $\{0,-1\}$
with the modular weights $\{0,1-k\}$ after that power forces
$k=2$ and power zero. Absorb the finite Dirichlet character into
the eigenform and take its primitive newform $f$. Li's Fricke formula
\cite[p.~296, before Lemma 6]{li-newforms} makes the span of its
coefficient conjugates Fricke-stable. Shimura's quotient construction
and Euler-factor identity \cite[Theorem 1]{shimura-factors} therefore
show that $r_5$ is a constituent of the Tate module of the modular abelian variety
$A_f$, or equivalently $r_5^\vee$ is its cohomological constituent.
This duality preserves the conductor.
